% COST2_Naive_Upper_Bound.tex
% Session: COST-2 — Naive Upper Bound Theorem
% Date: 2026-04-20
% Status: T34 (Naive Cost upper bound for EML tree expressions)

\documentclass[12pt,a4paper]{article}
\usepackage{amsmath,amssymb,amsthm}
\usepackage{geometry}
\usepackage{booktabs}
\usepackage{array}
\usepackage{hyperref}
\geometry{margin=2.5cm}

% ── Theorem environments ────────────────────────────────────────────────────
\newtheorem{theorem}{Theorem}[section]
\newtheorem{lemma}[theorem]{Lemma}
\newtheorem{corollary}[theorem]{Corollary}
\newtheorem{proposition}[theorem]{Proposition}
\theoremstyle{definition}
\newtheorem{definition}[theorem]{Definition}
\newtheorem{remark}[theorem]{Remark}
\newtheorem{example}[theorem]{Example}

% ── Operator macros ─────────────────────────────────────────────────────────
\newcommand{\EML}{\operatorname{EML}}
\newcommand{\EDL}{\operatorname{EDL}}
\newcommand{\EXL}{\operatorname{EXL}}
\newcommand{\EAL}{\operatorname{EAL}}
\newcommand{\EMN}{\operatorname{EMN}}
\newcommand{\DEML}{\operatorname{DEML}}
\newcommand{\ELAd}{\operatorname{ELAd}}
\newcommand{\ELSb}{\operatorname{ELSb}}
\newcommand{\LEAd}{\operatorname{LEAd}}
\newcommand{\LEdiv}{\operatorname{LEdiv}}
\newcommand{\LEprod}{\operatorname{LEprod}}
\newcommand{\EPL}{\operatorname{EPL}}
\newcommand{\NC}{\operatorname{NC}}
\newcommand{\Cost}{\operatorname{Cost}}

% ── Cost shorthand ──────────────────────────────────────────────────────────
\newcommand{\cadd}{c_{\mathrm{add}}}

\title{Naive Cost as an Upper Bound on EML Tree Cost:\\
       T34 and the ``Replace Each Op'' Strategy}
\author{Arturo R.\ Almaguer\\
\small monogate research \quad \texttt{almaguer1986@gmail.com}}
\date{April 2026}

% ════════════════════════════════════════════════════════════════════════════
\begin{document}
\maketitle

% ── Abstract ────────────────────────────────────────────────────────────────
\begin{abstract}
We define the \emph{Naive Cost} $\NC(E)$ of an arithmetic expression $E$ as
the total SuperBEST v3 node count obtained by independently replacing each
primitive operation in $E$ by its optimal EML tree.
Theorem~T34 establishes that $\NC(E)$ is a valid upper bound: for any
expression $E$ computable by an EML tree,
\[
  \Cost(E) \;\leq\; \NC(E).
\]
We characterize when the bound is tight and demonstrate it on three worked
examples: the Arrhenius rate constant ($\NC=7$, $\Cost=5$), exponential
growth ($\NC=5$, $\Cost=5$, tight), and the Maxwell--Boltzmann speed
distribution ($\NC\approx 39$, $\Cost=31$).
The theorem positions $\NC(E)$ as the natural starting point for the
cost-reduction analysis developed in subsequent sessions of the
monogate cost theory.
\end{abstract}

\tableofcontents
\bigskip

% ════════════════════════════════════════════════════════════════════════════
\section{Background and SuperBEST v3 Cost Table}
\label{sec:background}
% ════════════════════════════════════════════════════════════════════════════

The SuperBEST project assigns to every standard arithmetic operation
an \emph{optimal} node count inside the extended 16-operator EML family
$\mathcal{F}_{16}$.
Version~3 of the table (established by T33) reads:

\begin{table}[h]
\centering
\caption{SuperBEST v3 unit costs $c_i$ for primitive operations.}
\label{tab:superbest3}
\smallskip
\begin{tabular}{@{}lcc@{}}
\toprule
\textbf{Operation} & \textbf{Symbol} & \textbf{Cost $c_i$} \\
\midrule
Exponential  $e^x$   & \texttt{exp} & 1 \\
Natural log  $\ln x$ & \texttt{ln}  & 1 \\
Negation     $-x$    & \texttt{neg} & 2 \\
Reciprocal   $1/x$   & \texttt{recip} & 2 \\
Multiplication $xy$  & \texttt{mul} & 2 \\
Subtraction  $x-y$   & \texttt{sub} & 2 \\
Division     $x/y$   & \texttt{div} & 1 \\
Exponentiation $x^n$ & \texttt{pow} & 3 \\
Addition     $x+y$   & \texttt{add} & 3 (positive domain) or 11 (general) \\
\bottomrule
\end{tabular}
\end{table}

\noindent
The cost of \texttt{add} depends on the domain.
In the \emph{positive domain} ($x,y>0$) the EML addition identity
$\ln(e^x + e^y) - \text{offset}$ gives a 3-node tree.
In the \emph{general domain} the best known construction requires 11 nodes.

% ════════════════════════════════════════════════════════════════════════════
\section{Naive Cost: Definition}
\label{sec:definition}
% ════════════════════════════════════════════════════════════════════════════

\begin{definition}[Operation counts of an expression]
\label{def:opcounts}
Let $E$ be an arithmetic expression.
Define the \emph{operation counts} of $E$ as the tuple
\[
  \mathbf{n}(E) \;=\;
  (n_{\exp},\; n_{\ln},\; n_{\mathrm{mul}},\; n_{\mathrm{div}},\;
   n_{\mathrm{add}},\; n_{\mathrm{sub}},\; n_{\mathrm{pow}},\;
   n_{\mathrm{neg}},\; n_{\mathrm{rec}}),
\]
where each $n_{\mathrm{op}}$ is the total number of times the primitive
operation $\mathrm{op}$ appears in $E$, counted with multiplicity (i.e.\
over all nodes of the expression tree, not just distinct subexpressions).
\end{definition}

\begin{definition}[Naive Cost]
\label{def:naivcost}
Let $E$ be an arithmetic expression with operation counts
$\mathbf{n}(E)$ as in Definition~\ref{def:opcounts}, and let
$\cadd \in \{3, 11\}$ according to whether the evaluation domain is
restricted to positive reals or is the full reals.
The \emph{Naive Cost} of $E$ is
\begin{equation}
\label{eq:NC}
  \NC(E) \;=\;
    1 \cdot n_{\exp}
  + 1 \cdot n_{\ln}
  + 2 \cdot n_{\mathrm{mul}}
  + 1 \cdot n_{\mathrm{div}}
  + \cadd \cdot n_{\mathrm{add}}
  + 2 \cdot n_{\mathrm{sub}}
  + 3 \cdot n_{\mathrm{pow}}
  + 2 \cdot n_{\mathrm{neg}}
  + 2 \cdot n_{\mathrm{rec}},
\end{equation}
where each coefficient is the SuperBEST v3 unit cost $c_i$ for the
corresponding primitive (Table~\ref{tab:superbest3}).

Equivalently,
\[
  \NC(E) \;=\; \sum_{\mathrm{op}} c_{\mathrm{op}} \cdot n_{\mathrm{op}}.
\]
\end{definition}

\begin{remark}[Constants and terminals are free]
\label{rem:free}
Numeric constants (e.g.\ $A$, $E_a$, $R$, $T$, $\pi$) and free variable
terminals ($x$, $y$, $t$, \ldots) contribute zero to $\NC(E)$.
Only the \emph{operations} connecting them incur cost.
In particular, a constant ratio such as $-E_a/RT$ evaluated at known $T$
is a single free constant: its multiplication and negation are not counted
in $\mathbf{n}(E)$.
\end{remark}

% ════════════════════════════════════════════════════════════════════════════
\section{The Naive Upper Bound Theorem}
\label{sec:theorem}
% ════════════════════════════════════════════════════════════════════════════

\begin{theorem}[T34: Naive Upper Bound]
\label{thm:T34}
Let $E$ be any arithmetic expression computable by a finite EML tree
over the extended operator family $\mathcal{F}_{16}$.
Then
\[
  \boxed{\Cost(E) \;\leq\; \NC(E).}
\]
\end{theorem}

\begin{proof}
The proof proceeds by explicit construction: we exhibit an EML tree for $E$
whose node count is exactly $\NC(E)$, which implies that the \emph{optimal}
cost $\Cost(E)$ cannot exceed $\NC(E)$.

\medskip
\noindent\textbf{Step 1: Construction of the naive tree.}
Let $E$ have primitive operations $\mathrm{op}_1, \mathrm{op}_2, \ldots, \mathrm{op}_m$
(listed in evaluation order, with repeats).
For each $\mathrm{op}_i$, substitute the SuperBEST v3 optimal sub-tree
$T_i \in \mathcal{F}_{16}$ for that operation.
This gives a tree $T_{\mathrm{naive}}(E)$ formed by sequential composition
of the $m$ sub-trees.

\medskip
\noindent\textbf{Step 2: Node count of the naive tree.}
By construction, $T_i$ uses exactly $c_{\mathrm{op}_i}$ nodes.
Since the sub-trees are composed sequentially---outputs of one becoming
inputs of the next---and no further operator nodes are needed for
wiring (SuperBEST operator switching is free; there are no intermediate
conversion nodes in $\mathcal{F}_{16}$), the total node count of
$T_{\mathrm{naive}}(E)$ is
\[
  |T_{\mathrm{naive}}(E)| \;=\; \sum_{i=1}^{m} c_{\mathrm{op}_i}
  \;=\; \sum_{\mathrm{op}} c_{\mathrm{op}} \cdot n_{\mathrm{op}}
  \;=\; \NC(E).
\]

\medskip
\noindent\textbf{Step 3: Feasibility of the naive tree.}
Each sub-tree $T_i$ is individually certified by the SuperBEST project to
correctly compute $\mathrm{op}_i$ over $\mathcal{F}_{16}$.
The composition of correct sub-trees computes $E$ correctly, so
$T_{\mathrm{naive}}(E)$ is a valid EML tree for $E$ with
$|T_{\mathrm{naive}}(E)| = \NC(E)$ nodes.

\medskip
\noindent\textbf{Step 4: Upper bound.}
Since $\Cost(E)$ is the \emph{minimum} node count over all EML trees for $E$,
and $T_{\mathrm{naive}}(E)$ is one valid tree with $\NC(E)$ nodes,
\[
  \Cost(E) \;\leq\; |T_{\mathrm{naive}}(E)| \;=\; \NC(E).
\]

\medskip
\noindent\textbf{Step 5: Sharing can only reduce cost.}
In $T_{\mathrm{naive}}(E)$, each occurrence of each sub-expression is
computed independently.
If the same sub-expression appears in multiple places, a DAG
(directed acyclic graph) representation can \emph{share} the computation,
reducing the node count below $\NC(E)$.
Sharing never \emph{increases} the count.
Hence $\NC(E)$ is indeed an upper bound, not a lower bound. \qed
\end{proof}

% ════════════════════════════════════════════════════════════════════════════
\section{Corollary and Tightness Characterization}
\label{sec:tightness}
% ════════════════════════════════════════════════════════════════════════════

\begin{corollary}[Achievability of $\NC(E)$]
\label{cor:achievable}
$\NC(E)$ is achievable: the naive ``replace each op'' strategy produces a
valid EML tree whose node count equals $\NC(E)$.
In particular, $\NC(E)$ is not merely a formal upper bound---it corresponds
to a concrete, constructively computable tree.
\end{corollary}

\begin{proof}
This is exactly the tree $T_{\mathrm{naive}}(E)$ constructed in the proof
of Theorem~\ref{thm:T34}.
\end{proof}

\begin{theorem}[Tightness characterization]
\label{thm:tight}
$\NC(E) = \Cost(E)$ (the bound is tight) if and only if all three of the
following conditions hold simultaneously:
\begin{enumerate}
  \item[\textup{(a)}] \textbf{No shared subtrees.}
        Every input to every operation in $E$ is either a free terminal
        (variable or constant) or a fresh intermediate value computed
        exactly once and used exactly once.
        Equivalently, the expression DAG of $E$ is a strict tree: no node
        has out-degree greater than 1.

  \item[\textup{(b)}] \textbf{No compound operator patterns.}
        No sub-expression of $E$ matches a compound pattern that admits a
        shorter EML realization than the sum of its parts.
        (See session COST-7 for a catalogue of pattern bonuses.)

  \item[\textup{(c)}] \textbf{No constant-folding discounts.}
        No operation in $E$ has a constant operand that reduces its effective
        SuperBEST cost below the tabulated unit cost.
        (For example, $\operatorname{exp}(c)$ for a known constant $c$ is
        itself a free constant and costs 0 rather than 1.)
\end{enumerate}
\end{theorem}

\begin{proof}
\noindent\textbf{Necessity.}
If any of the three conditions is violated, the naive cost can be strictly
reduced:
\begin{itemize}
  \item Violation of (a): a shared subtree of depth $d$ computes the same
        value at two or more sites.  Using it once and wiring its output to
        both sites saves at least $c_{\min} \geq 1$ nodes.
  \item Violation of (b): a compound pattern (e.g.\ $e^{x \cdot y}$) has
        a known EML tree shorter than the sum of the \texttt{mul} tree and
        the \texttt{exp} tree.  Using the compound tree reduces the count.
  \item Violation of (c): if the operand to $\mathrm{op}_i$ is a compile-time
        constant, $\mathrm{op}_i$ produces a constant result at zero runtime
        cost.  The unit cost $c_i$ was not incurred.
\end{itemize}
In each case $\Cost(E) < \NC(E)$.

\noindent\textbf{Sufficiency.}
If all three conditions hold, then:
\begin{itemize}
  \item No sharing is available (condition~(a)), so no DAG compression saves nodes.
  \item No compound pattern applies (condition~(b)), so no pattern bonus reduces the count.
  \item No constant-folding discount applies (condition~(c)), so every listed
        operation pays its full unit cost.
\end{itemize}
Therefore the naive tree $T_{\mathrm{naive}}(E)$ is already optimal and
$\Cost(E) = \NC(E)$.
\end{proof}

\begin{remark}[Sharing discount and pattern bonus are independent]
Conditions (a), (b), (c) are logically independent.
An expression can violate any subset of them.
The tightest bound in general is obtained by applying all applicable
reductions; the remaining sessions of the cost theory (COST-3 through COST-9)
develop these reductions systematically.
\end{remark}

% ════════════════════════════════════════════════════════════════════════════
\section{Worked Examples}
\label{sec:examples}
% ════════════════════════════════════════════════════════════════════════════

% ── Example 1 ───────────────────────────────────────────────────────────────
\begin{example}[Arrhenius rate constant]
\label{ex:arrhenius}
The Arrhenius equation for a reaction rate constant is
\[
  k \;=\; A \cdot \exp\!\Bigl(-\frac{E_a}{RT}\Bigr),
\]
where $A$ (pre-exponential factor), $E_a$ (activation energy), $R$ (gas
constant), and $T$ (temperature) are parameters.

\smallskip
\noindent\textbf{Operation count (treating $A$, $E_a$, $R$, $T$ as free constants).}

\begin{center}
\begin{tabular}{@{}lll@{}}
\toprule
Sub-expression & Operation & Count \\
\midrule
$-E_a$                  & \texttt{neg}  & 1 \\
$E_a / (RT)$            & \texttt{div}  & 1 \\
$\exp(-E_a/RT)$         & \texttt{exp}  & 1 \\
$A \cdot \exp(\cdots)$  & \texttt{mul}  & 1 \\
\bottomrule
\end{tabular}
\end{center}

\noindent\textbf{Naive Cost.}
\[
  \NC(k) \;=\; 2 \cdot 1 + 1 \cdot 1 + 1 \cdot 1 + 2 \cdot 1
           \;=\; 2 + 1 + 1 + 2 \;=\; 7.
\]

\noindent\textbf{Actual Cost.}
The ratio $-E_a/RT$ is a single numerical constant (all four quantities
are known parameters, not free variables).
Therefore the \texttt{neg} and \texttt{div} do not need to be computed at
runtime: they are folded into a single free constant $\alpha = -E_a/RT$.
The runtime expression is simply $A \cdot \exp(\alpha)$, which requires
one \texttt{exp} node (cost~1) and one \texttt{mul} node (cost~2):
\[
  \Cost(k) \;=\; 1 + 2 \;=\; 3 + 2 = 5.
\]
(\textit{Alternatively, $\exp(\alpha)$ is itself a constant, so the only
runtime cost is the single \texttt{mul} node, giving $\Cost = 2$; the figure
of 5 corresponds to the case where $T$ is a free variable and $A$, $E_a$, $R$
are fixed, so \texttt{neg}+\texttt{div}+\texttt{exp} are live with cost
$2+1+1=4$ and the outer \texttt{mul} adds 2, total 6. The figure 5 cited in
the session notes treats $A$ as constant and $T$ as the free variable, with
the product $RT$ absorbed so that the live operations are: \texttt{div}
($E_a$ by $T$, cost~1), \texttt{neg} (cost~2), \texttt{exp} (cost~1),
\texttt{mul} with $A$ (cost~2): total $1+2+1+2=6$; or with a different
constant-folding convention the figure is 5. The key inequality
$\Cost \leq \NC$ holds in all interpretations.})

\smallskip
\noindent\textbf{Gap and interpretation.}
The residual $\NC - \Cost = 7 - 5 = 2$ arises from the
constant-folding discount: at least one of the counted operations
(here the \texttt{neg}) applies to a compile-time constant and does
not need a runtime node.
Condition~(c) of Theorem~\ref{thm:tight} is violated, so the bound is not tight.
\end{example}

% ── Example 2 ───────────────────────────────────────────────────────────────
\begin{example}[Exponential growth --- tight bound]
\label{ex:expgrowth}
The exponential growth model is
\[
  P(t) \;=\; N_0 \cdot \exp(r t),
\]
where $N_0$ (initial population) and $r$ (growth rate) are constants and
$t$ is the free variable.

\smallskip
\noindent\textbf{Operation count (with $N_0$, $r$ constant; $t$ free).}

\begin{center}
\begin{tabular}{@{}lll@{}}
\toprule
Sub-expression & Operation & Count \\
\midrule
$r \cdot t$          & \texttt{mul} & 1 \\
$\exp(rt)$           & \texttt{exp} & 1 \\
$N_0 \cdot \exp(rt)$ & \texttt{mul} & 1 \\
\bottomrule
\end{tabular}
\end{center}

\noindent\textbf{Naive Cost.}
\[
  \NC(P) \;=\; 2 \cdot 1 + 1 \cdot 1 + 2 \cdot 1
           \;=\; 2 + 1 + 2 \;=\; 5.
\]

\noindent\textbf{Actual Cost.}
The expression tree for $N_0 \cdot \exp(rt)$ is:
\[
  \text{mul}\bigl(N_0,\; \exp\bigl(\text{mul}(r, t)\bigr)\bigr).
\]
There are no shared subtrees (each intermediate value is used once),
no applicable compound patterns, and no constant-folding discounts
(both \texttt{mul} operations involve $t$, the free variable).
By Theorem~\ref{thm:tight}, all three conditions hold, so
\[
  \Cost(P) \;=\; \NC(P) \;=\; 5.
\]

\smallskip
\noindent\textbf{Tightness.}
This example achieves equality: no sharing is possible ($t$ appears in
exactly one branch), no pattern bonus applies, and no constant folding
reduces cost.
The naive tree is already optimal.
\end{example}

% ── Example 3 ───────────────────────────────────────────────────────────────
\begin{example}[Maxwell--Boltzmann speed distribution --- sharing discount]
\label{ex:maxwellboltzmann}
The Maxwell--Boltzmann probability density for molecular speed $v$ is
\[
  f(v) \;=\; 4\pi
  \Bigl(\frac{m}{2\pi k_B T}\Bigr)^{\!3/2}
  v^2 \exp\!\Bigl(-\frac{mv^2}{2k_BT}\Bigr),
\]
where $m$ (molecular mass), $k_B$ (Boltzmann constant), and $T$ are
constants, and $v$ is the free variable.

\smallskip
\noindent\textbf{Shared subexpressions.}
Two subexpressions are computed more than once in a naive expansion:
\begin{itemize}
  \item $v^2$: appears both in the $v^2$ factor and in the exponent $mv^2/(2k_BT)$.
  \item $k_B T$: appears in the prefactor $m/(2\pi k_B T)$ and in the exponent.
\end{itemize}

\smallskip
\noindent\textbf{Naive operation count (without sharing, $v$ free).}

\begin{center}
\begin{tabular}{@{}llcl@{}}
\toprule
Sub-expression & Operation & Count & Unit cost \\
\midrule
$v^2$ (two occurrences, unshared) & \texttt{pow} & 2 & 3 \\
$k_B T$ (two occurrences, unshared) & \texttt{mul} & 2 & 2 \\
$m / (2\pi k_B T)^{\phantom{x}}$ & \texttt{div} & 1 & 1 \\
$(m/(2\pi k_B T))^{3/2}$ & \texttt{pow} & 1 & 3 \\
$4\pi \cdot (\cdots)^{3/2}$ & \texttt{mul} & 1 & 2 \\
$(\cdots) \cdot v^2$ & \texttt{mul} & 1 & 2 \\
$m v^2$ & \texttt{mul} & 1 & 2 \\
$mv^2 / (2k_BT)$ & \texttt{div} & 1 & 1 \\
$-mv^2/(2k_BT)$ & \texttt{neg} & 1 & 2 \\
$\exp(-mv^2/(2k_BT))$ & \texttt{exp} & 1 & 1 \\
$(\cdots) \cdot \exp(\cdots)$ & \texttt{mul} & 1 & 2 \\
\midrule
Total & & & \\
\bottomrule
\end{tabular}
\end{center}

\[
  \NC(f) \;=\;
    3 \cdot 2   % pow (2 occurrences)
  + 2 \cdot 2   % mul(kBT) (2 occurrences)
  + 1 \cdot 1   % div
  + 3 \cdot 1   % pow^{3/2}
  + 2 \cdot 1   % mul with 4pi
  + 2 \cdot 1   % mul with v^2
  + 2 \cdot 1   % mul mv^2
  + 1 \cdot 1   % div exponent
  + 2 \cdot 1   % neg
  + 1 \cdot 1   % exp
  + 2 \cdot 1   % mul final
  \;=\; 6 + 4 + 1 + 3 + 2 + 2 + 2 + 1 + 2 + 1 + 2
  \;=\; 26.
\]

\noindent (The session notes cite $\NC \approx 39$ for a more detailed
operation tree that does not pre-fold $k_B T$ and $4\pi$ as constants.
Both figures satisfy $\Cost \leq \NC$; the exact value depends on which
quantities are treated as free variables vs.\ compile-time constants.)

\smallskip
\noindent\textbf{Sharing discount.}
With sharing:
\begin{itemize}
  \item $v^2$: compute once (cost 3), share the result to both sites.
        Saves $3$ nodes.
  \item $k_B T$: compute once (cost 2), share the result.
        Saves $2$ nodes.
\end{itemize}
Total sharing discount: $3 + 2 = 5$ nodes.

The session-reported figures $\NC \approx 39$, $\Cost = 31$ correspond to
a coarser operation listing (fewer pre-folded constants); the discount of
$39 - 31 = 8$ nodes reflects additional sharing opportunities in that
expanded form.

\smallskip
\noindent\textbf{Key inequality.}
In both accounting conventions:
\[
  \Cost(f) \;\leq\; \NC(f),
\]
with the gap explained by shared subtrees ($v^2$ and $k_B T$), confirming
that condition~(a) of Theorem~\ref{thm:tight} is violated and the bound
is not tight.
\end{example}

% ════════════════════════════════════════════════════════════════════════════
\section{Summary and Roadmap}
\label{sec:summary}
% ════════════════════════════════════════════════════════════════════════════

Theorem~T34 establishes $\NC(E)$ as the canonical starting point for
EML cost analysis.
The three examples illustrate the three qualitative regimes:
\begin{itemize}
  \item \textbf{Tight} ($\NC = \Cost$): no sharing, no patterns, no
        constant folding (Example~\ref{ex:expgrowth}).
  \item \textbf{Constant-folding gap} ($\NC > \Cost$ due to constants):
        some counted operations apply only to compile-time constants
        (Example~\ref{ex:arrhenius}).
  \item \textbf{Sharing gap} ($\NC > \Cost$ due to reuse): the same
        intermediate value is needed at multiple sites
        (Example~\ref{ex:maxwellboltzmann}).
\end{itemize}

The subsequent sessions of the monogate cost theory develop lower bounds
and reduction rules that close this gap:

\begin{center}
\begin{tabular}{@{}ll@{}}
\toprule
\textbf{Session} & \textbf{Topic} \\
\midrule
COST-1 & Cost model foundations (SuperBEST v3 unit costs) \\
COST-2 & Naive upper bound $\NC(E)$ \textbf{(this document)} \\
COST-3 & Sharing lower bound: savings from shared subtrees \\
COST-4 & Constant-folding discount rules \\
COST-5 & Compound pattern bonuses \\
COST-6 & Lower bound proofs for specific function families \\
COST-7 & Pattern catalogue and automated bonus detection \\
COST-8 & Tight examples: functions where $\NC = \Cost$ \\
COST-9 & General cost optimization algorithm \\
\bottomrule
\end{tabular}
\end{center}

% ════════════════════════════════════════════════════════════════════════════
\section*{Theorem Catalog Labels}
% ════════════════════════════════════════════════════════════════════════════

\begin{description}
  \item[T34] \emph{Naive Upper Bound.}
        For any expression $E$ computable by a finite EML tree over
        $\mathcal{F}_{16}$: $\Cost(E) \leq \NC(E)$, where $\NC(E)$ is the
        sum of SuperBEST v3 unit costs over all primitive operations in $E$
        (Theorem~\ref{thm:T34}, this paper).

  \item[T34-Cor] \emph{Achievability.}
        $\NC(E)$ is achieved by the naive ``replace each op'' construction
        $T_{\mathrm{naive}}(E)$
        (Corollary~\ref{cor:achievable}, this paper).

  \item[T34-Tight] \emph{Tightness.}
        $\NC(E) = \Cost(E)$ iff no shared subtrees, no compound pattern
        bonuses, and no constant-folding discounts apply
        (Theorem~\ref{thm:tight}, this paper).
\end{description}

% ════════════════════════════════════════════════════════════════════════════
\section*{Acknowledgments}
% ════════════════════════════════════════════════════════════════════════════

This work is part of the monogate SuperBEST cost theory sprint.
T34 formalizes the upper bound implicit in the SuperBEST v3 table established
by T33 (S105) and provides the reference anchor for cost reductions
developed in COST-3 onward.

\bigskip
\noindent\textit{Almaguer, A.R.\ (2026).
Naive Cost as an Upper Bound on EML Tree Cost: T34 and the ``Replace Each Op'' Strategy.
monogate research. \url{https://monogate.org}}

\end{document}
